%% TRANSCRIBED (PDF 87-91)
\apphead{Killing arithmetic groups in tori}{KILLING ARITHMETIC GROUPS IN TORI}

\subsechead{A.1}{Arithmetic groups in tori}

Let $A$ be a linear algebraic group over $\Q$, and let $\Gamma$ be a subgroup
of the group $A(\Q)$ of rational points of $A$. Then $\Gamma$ is said to be an
\emph{arithmetic subgroup} if for any algebraic embedding $A\subset\GL_n$
($n$ arbitrary) the groups $\Gamma$ and $A(\Q)\cap\GL_n(\Z)$ are
\emph{commensurable} (two subgroups $\Gamma_1,\Gamma_2$ are said to be
commensurable if $\Gamma_1\cap\Gamma_2$ is of finite index in $\Gamma_1$ and
$\Gamma_2$). It is well-known that it suffices to check that $\Gamma$ and
$A(\Q)\cap\GL_n(\Z)$ are commensurable for one embedding $A\subset\GL_n$.

\runhead{Example}

Let $K$ be a number field and let $E$ be the group of units of $K$. Then $E$
is an arithmetic subgroup of $T=\mathrm{R}_{K/\Q}(\Gm)$.

If $T$ is a torus over $\Q$, let $T^o$ be the intersection of the kernels of
the homomorphisms of $T$ into $\Gm$. The torus $T$ is said to be
\emph{anisotropic} if $T=T^o$; in terms of the character group $X=X(T)$ this
means that $X$ has no non-zero elements which are left fixed by
$G=\Gal(\bar{\Q}/\Q)$.

\begin{theorem}
Let $T$ be a torus over $\Q$, and let $\Gamma$ be an arithmetic subgroup of
$T$. Then $\Gamma\cap T^o$ is of finite index in $\Gamma$, and the quotient
$T^o(R)/\Gamma\cap T^o$ is compact.
\end{theorem}

This is due to T.~Ono; for a proof of a more general statement
(``Godement's conjecture'') see Mostow--Tamagawa \cite{mostow-tamagawa}.

\begin{corollary}
Let $T$ be a torus over $\Q$, and let $\Gamma$ be an arithmetic subgroup of
$T$. If $T$ is anisotropic, then $T(R)/\Gamma$ is compact.
\end{corollary}

\runhead{Exercise}

Let $T$ be a torus over $\Q$, with character group $X$.
\begin{alist}
\item[\emph{a})] Show that
\[
T(\Q)=\Hom_{\Gal}(X,\bar{\Q}^*).
\]
\item[\emph{b})] Let $U$ be the subgroup of $\bar{\Q}^*$ whose elements are
the algebraic units of $\bar{\Q}$. Let
\[
\Gamma=\Hom_{\Gal}(X,U).
\]
Show that $\Gamma$ is an arithmetic subgroup of $T(\Q)$ and that any
arithmetic subgroup of $T(\Q)$ is contained in $\Gamma$.
\end{alist}

\subsechead{A.2}{Killing arithmetic subgroups}

Let $T$ be a torus over $\Q$, and let $X(T)$ be its character group; put
$Y(T)=X(T)\tens_\Z\Q$. Let $\Lambda$ be the set of classes of
$\Q$-irreducible representations of $G=\Gal(\bar{\Q}/\Q)$ through its finite
quotients. For each $\lambda\in\Lambda$, let $Y_\lambda$ be the corresponding
isotypic sub-$G$-module of $Y$, i.e. the sum of all sub-$G$-modules of $Y$
isomorphic to $\lambda$. One has the direct sum decomposition
\[
Y=\coprod_{\lambda\in\Lambda} Y_\lambda.
\]

Let $Y^o=Y_1$, where $1$ is the unit representation of $G$; let $Y^-$ be the
sum of those $Y_\lambda$ where for all the infinite Frobeniuses
$c\in C_\infty$ (cf.~\S\ref{sec:II.3.1}) we have $\lambda(c)=-1$; let $Y^+$
be the sum of the other $Y_\lambda$. We have
\begin{align*}
&Y^o=Y^G=\{y\in Y \mid gy=y \quad\text{for all}\quad g\in G\}\\
&Y^-=\{y\in Y \mid cy=-y \qquad\text{for all}\quad c\in C_\infty\},\\
&Y=Y^o\oplus Y^-\oplus Y^+.
\end{align*}

Note that $Y=Y^o$ if and only if $T$ is anisotropic.

If $c\in C_\infty$, and $H=\{1,c\}$, then, since
$T(R)=\Hom_H(X(T),C^*)$, we see that $T(R)$ \emph{is compact if and only if}
$Y=Y^-$.

\begin{proposition}
Let $\Gamma$ be an arithmetic subgroup of the torus $T$, and $\bar{\Gamma}$
its Zariski closure (cf.~\S\ref{sec:II.1.2}). Then:
\begin{equation}
\tag{$*$}\label{eq:II.A.2.star}
Y(T/\bar{\Gamma})=Y^o\oplus Y^-.
\end{equation}

{\upshape[Since the torus $T/\bar{\Gamma}$ is a quotient of $T$, we identify
$Y(T/\bar{\Gamma})$ with a submodule of $Y(T)$.]}
\end{proposition}

\begin{proof}
Suppose first that $Y$ is \emph{irreducible}, i.e. that $T$ has no proper
subtori and is $\neq 0$.

If $Y=Y^o$, then $T$ is isomorphic to $\Gm$ and hence $\Gamma$ is finite.
This shows that $Y(T/\bar{\Gamma})=Y(T)$, hence \eqref{eq:II.A.2.star}.
If $Y=Y^-$, then $T(R)$ is compact. Since $\Gamma$ is a discrete subgroup of
$T(R)$, it is finite. Hence $Y(T/\bar{\Gamma})=Y(T)$ and
\eqref{eq:II.A.2.star} follows.

If $Y=Y^+$, then $T(R)$ is not compact. Consequently, $\Gamma$ is infinite
since $T(R)/\Gamma$ is compact by Ono's theorem. Hence $\bar{\Gamma}$ is an
algebraic subgroup of $T$ of dimension $\geq 1$. Its connected component is
a non-trivial subtorus of $T$. This shows that $\bar{\Gamma}=T$, hence
$Y(T/\bar{\Gamma})=0$. Hence again \eqref{eq:II.A.2.star}.

The general case follows easily from the irreducible one; for instance,
choose a torus $T'$ to $T$ which splits in direct product of irreducible
tori and note that $\Gamma$ is commensurable with the image by $T'\to T$ of
an arithmetic subgroup of $T$.
%% p91: "T' to T" above is as printed (the word "to", verified at dpi=500;
%% the arrow glyph $T'\to T$ occurs only in the following sentence).
\end{proof}

\runhead{Exercise}

Let $y\in Y$. Define $Ny$ as the mean value of the transforms of $y$ by $G$.
\begin{alist}
\item[\emph{a}] Prove that $N$ is a $G$-linear projection of $Y$ onto $Y^o$,
hence $\Ker(N)=Y^-\oplus Y^+$.
\item[\emph{b}] Prove that $Y^+$ is generated by the elements $cy+y$, with
$y\in\Ker(N)$, $c\in C_\infty$.
\end{alist}
